\section{Limitações}

A avaliação da qualidade em uso por meio de métricas quantitativas apresenta desafios, dada a sua natureza subjetiva (como apontado por \citeonline{qatch} e \citeonline{fabijan_three_2019}). Embora as métricas adotadas sejam baseadas em modelos de referência, não foi possível abranger todas as características subcaracterísticas de qualidade em uso, o que seria ideal para uma avaliação mais abrangente e assertiva.

Quanto ao suporte à tomada de decisão, a principal limitação foi o tempo disponível para esta pesquisa. Os dados quantitativos indicaram a necessidade de uma análise mais aprofundada da funcionalidade, sugerindo que novas iterações poderiam ser conduzidas com insumos qualitativos da opinião dos usuários. Avaliações futuras são essenciais para compreender melhor as diferenças de comportamento dos usuários e suas implicações no contexto do negócio.

\section{Trabalhos Futuros}

Devido às restrições do cronograma deste trabalho de conclusão de curso, algumas etapas do processo proposto não puderam ser exploradas em profundidade. Entre elas, destaca-se a coleta de dados qualitativos sobre as versões do produto avaliadas, o que poderia fornecer informações adicionais para embasar a tomada de decisão (conforme ilustrado na Figura \ref{fig:processo-novo}).

Além disso, todo o processo de análise e decisão foi conduzido manualmente, em função das limitações das ferramentas utilizadas e do volume de dados coletado. Estudos futuros podem explorar a automação dos procedimentos de análise, além da liberação da versão definitiva, se utilizando de modelos de aprendizado de máquina integrados a rotinas de entrega contínua.
